\documentclass[a4paper, 12pt]{article}

% --- Pacotes Essenciais ---
\usepackage[utf8]{inputenc}
\usepackage[portuguese]{babel}
\usepackage{geometry}
\usepackage{amsmath} % Pacote principal para matemática
\usepackage{amssymb} % Para símbolos como \mathbb{R}

% --- Configuração de Página ---
\geometry{a4paper, margin=1in}

% --- Informações do Documento ---
\title{Derivação da Lei de Controle Cinemático \\
       para um Manipulador 2-DOF \\
       \large (Mínimos Quadrados Amortecidos - DLS)}
\author{Mendes, J}
\date{\today}

% --- Início do Documento ---
\begin{document}

\maketitle

\section{Introdução}
O objetivo é controlar a posição do efetuador final (EE) de um manipulador robótico de 2 graus de liberdade (2-DOF) para que ele alcance uma posição de referência $\mathbf{x}_d$ no espaço cartesiano (espaço de tarefa).

A dificuldade central é que controlamos diretamente as variáveis de junta, $\mathbf{q} \in \mathbb{R}^2$, e não a posição cartesiana $\mathbf{x} \in \mathbb{R}^2$. A lei de controle deve, portanto, "traduzir" um erro no espaço de tarefa em comandos de velocidade para o espaço de juntas.

Utilizamos uma abordagem de controle em malha fechada no espaço de tarefa, onde a solução da cinemática inversa é calculada iterativamente usando o método de Mínimos Quadrados Amortecidos (Damped Least Squares - DLS), que é robusto a singularidades.

\section{Modelo Cinemático}

\subsection{Cinemática Direta (Posição)}
A relação não linear entre as posições das juntas $\mathbf{q}$ e a posição cartesiana do efetuador final $\mathbf{x}$ é dada pela função de cinemática direta, $f$:
\[
\mathbf{x} = f(\mathbf{q})
\]
Onde $\mathbf{q} = \begin{bmatrix} q_1 \\ q_2 \end{bmatrix}$ e $\mathbf{x} = \begin{bmatrix} x \\ y \end{bmatrix}$.

\subsection{Cinemática Diferencial (Velocidade)}
Para relacionar as velocidades, derivamos a equação acima em relação ao tempo:
\[
\dot{\mathbf{x}} = \frac{d}{dt} f(\mathbf{q}) = \frac{\partial f(\mathbf{q})}{\partial \mathbf{q}} \frac{d\mathbf{q}}{dt}
\]
Isso nos dá a relação fundamental da cinemática diferencial, que envolve a matriz Jacobiana $\mathbf{J}(\mathbf{q})$:
\begin{equation}
\dot{\mathbf{x}} = \mathbf{J}(\mathbf{q}) \dot{\mathbf{q}}
\label{eq:jacobiana}
\end{equation}
Onde $\mathbf{J}(\mathbf{q}) \in \mathbb{R}^{2 \times 2}$ é a Jacobiana, que mapeia velocidades de junta ($\dot{\mathbf{q}}$) para velocidades cartesianas ($\dot{\mathbf{x}}$).

\section{Formulação da Lei de Controle}
A estratégia de controle é dividida em duas etapas:
\begin{enumerate}
    \item Definir uma velocidade cartesiana desejada ($\dot{\mathbf{x}}_d$) que leve o efetuador em direção ao alvo.
    \item Calcular as velocidades de junta necessárias ($\dot{\mathbf{q}}$) para realizar essa velocidade cartesiana.
\end{enumerate}

\subsection{Etapa 1: Controle Proporcional no Espaço de Tarefa}
Primeiro, definimos o erro de posição no espaço de tarefa como a diferença vetorial entre a posição desejada (alvo) $\mathbf{x}_d$ e a posição atual $\mathbf{x}$:
\begin{equation}
\mathbf{e} = \mathbf{x}_d - \mathbf{x}
\label{eq:erro}
\end{equation}
Para fazer o robô se mover em direção ao alvo, definimos uma velocidade desejada para o efetuador, $\dot{\mathbf{x}}_d$, que seja proporcional ao erro. Esta é a essência de um controlador Proporcional (P):
\begin{equation}
\dot{\mathbf{x}}_d = K_p \mathbf{e}
\label{eq:p_control}
\end{equation}
Onde $K_p$ é um ganho escalar positivo. Substituindo \eqref{eq:erro}, temos:
\[
\dot{\mathbf{x}}_d = K_p (\mathbf{x}_d - \mathbf{x})
\]
Isso cria um sistema de primeira ordem que, se $\dot{\mathbf{x}} = \dot{\mathbf{x}}_d$, garante que o erro $\mathbf{e}(t)$ convirja exponencialmente para zero.

\subsection{Etapa 2: Cinemática Inversa com DLS}
Agora, precisamos encontrar o $\dot{\mathbf{q}}$ que satisfaça a equação \eqref{eq:jacobiana} usando a velocidade desejada $\dot{\mathbf{x}}_d$ que acabamos de definir:
\[
\mathbf{J}(\mathbf{q}) \dot{\mathbf{q}} = \dot{\mathbf{x}}_d
\]
O desafio é inverter $\mathbf{J}(\mathbf{q})$. Uma inversão simples, $\dot{\mathbf{q}} = \mathbf{J}^{-1} \dot{\mathbf{x}}_d$, falha catastroficamente quando o robô está em uma configuração singular (por exemplo, braço totalmente esticado), pois o determinante de $\mathbf{J}$ vai a zero e $\mathbf{J}^{-1}$ "explode".

Para evitar isso, usamos uma solução de otimização mais robusta. Em vez de resolver a equação exatamente, minimizamos uma função de custo que equilibra dois objetivos:
\begin{enumerate}
    \item Minimizar o erro da velocidade: $|| \mathbf{J}\dot{\mathbf{q}} - \dot{\mathbf{x}}_d ||^2$
    \item Minimizar o "esforço" (magnitude das velocidades das juntas): $||\dot{\mathbf{q}}||^2$
\end{enumerate}
A função de custo combinada é:
\begin{equation}
L(\dot{\mathbf{q}}) = || \mathbf{J}\dot{\mathbf{q}} - \dot{\mathbf{x}}_d ||^2 + \lambda^2 ||\dot{\mathbf{q}}||^2
\label{eq:custo}
\end{equation}
Onde $\lambda$ (lambda) é um fator de amortecimento (ou regularização). Este é o problema de Mínimos Quadrados Amortecidos (Damped Least Squares - DLS).

A solução analítica que minimiza \eqref{eq:custo} é dada por:
\begin{equation}
\dot{\mathbf{q}} = \mathbf{J}^T (\mathbf{J} \mathbf{J}^T + \lambda^2 \mathbf{I})^{-1} \dot{\mathbf{x}}_d
\label{eq:dls_sol}
\end{equation}
Onde $\mathbf{J}^T$ é a transposta da Jacobiana e $\mathbf{I}$ é a matriz identidade $2 \times 2$.

\textbf{Por que isso é robusto?} O termo $\lambda^2 \mathbf{I}$ é adicionado à matriz $\mathbf{J} \mathbf{J}^T$. Mesmo que $\mathbf{J} \mathbf{J}^T$ se torne singular (perdendo posto), a matriz $(\mathbf{J} \mathbf{J}^T + \lambda^2 \mathbf{I})$ é sempre positiva definida e, portanto, sempre invertível. Isso "amortece" a solução e evita que $\dot{\mathbf{q}}$ assuma valores infinitos.

\section{Implementação Discreta na Simulação}
Na simulação, o controle não é contínuo, mas executado em pequenos passos de tempo $\Delta t$. A lei de controle é implementada como um algoritmo iterativo:

\begin{enumerate}
    \item \textbf{Obter Posições:} \\
    Posição desejada: $\mathbf{x}_d$ (definida pelo clique do mouse) \\
    Posição atual: $\mathbf{x}_k = f(\mathbf{q}_k)$ (calculada pela cinemática direta)
    
    \item \textbf{Calcular Erro e Velocidade Desejada:} \\
    $\mathbf{e}_k = \mathbf{x}_d - \mathbf{x}_k$ \\
    $\dot{\mathbf{x}}_d = K_p \mathbf{e}_k$
    
    \item \textbf{Verificar Parada:} \\
    Se $|| \mathbf{e}_k || < \text{TOLERANCE}$, o alvo foi alcançado. Parar.
    
    \item \textbf{Calcular Jacobiana:} \\
    $\mathbf{J}_k = \mathbf{J}(\mathbf{q}_k)$
    
    \item \textbf{Resolver DLS para Velocidades de Junta:} \\
    $\dot{\mathbf{q}}_k = \mathbf{J}_k^T (\mathbf{J}_k \mathbf{J}_k^T + \lambda^2 \mathbf{I})^{-1} \dot{\mathbf{x}}_d$
    
    \item \textbf{Integrar (Método de Euler) para o Próximo Estado:} \\
    $\mathbf{q}_{k+1} = \mathbf{q}_k + \dot{\mathbf{q}}_k \cdot \Delta t$
    
    \item \textbf{Repetir} no próximo ciclo de simulação.
\end{enumerate}

\section{Conclusão}
A lei de controle implementada é um controlador P no espaço de tarefa, mapeado para o espaço de juntas através da cinemática inversa diferencial. A robustez do sistema é garantida pelo uso da formulação DLS (Mínimos Quadrados Amortecidos) para o cálculo da inversa da Jacobiana, o que evita instabilidades em configurações singulares e produz o movimento suave observado na simulação.

\end{document}